\section{Impossibility of global pairing}
\label{sec:no-pairing}

A pairing strategy preassigns disjoint pairs and answers each attacker stone
in a paired cell by taking its mate.  In a single-placement game this gives a
perfect local defense.  Under Hexo's two-placement turns, the attacker may
occupy both cells of a pair within one turn, so even an existing pairing would
need a separate argument.  The theorem below makes the question moot: the
required matching does not exist.

\begin{theorem}[No pairing; source T8]
\label{thm:T8}
There is no partial matching $M$ on the cells, with each cell in at most one
pair, such that every window contains both cells of some pair.
\end{theorem}

\begin{proof}
Discard every pair that is not contained in a window.  This removes every
non-axis-aligned pair and every axis pair at distance greater than $5$, while
preserving the covering property.  By \cref{fact:F1}, each remaining pair is
axis-aligned at an axis-distance $\delta\in\{1,\ldots,5\}$.

Fix a lattice line $\ell$.  A pair at positions $(p,p+\delta)$ is contained
in the windows whose integer starts satisfy
\[
  s\in[p+\delta-5,p].
\]
There are exactly $6-\delta\le5$ such starts.  If the left cells of
consecutive pairs are $p_i$ and $p_{i+1}$, coverage of every start requires
$p_{i+1}-p_i\le6-\delta_{i+1}\le5$.  Allowing mixed values of $\delta$ cannot
improve this bound.  A segment of $\ell$ with $L$ cells has $L-5$ internal
window starts and therefore contains at least $(L-5)/5$ pairs.  At least
$2(L-5)/5$ of its cells are matched along the axis of $\ell$.

Each matched cell serves exactly one axis, namely the axis of its pair.  Let
\[
  B_R=\{x:\dist(x,0)\le R\},
  \qquad |B_R|=3R^2+3R+1.
\]
For each of the three axes, $B_R$ meets exactly $2R+1$ line segments, and the
lengths of those segments sum to $|B_R|$.  Summing the line bound over all
segments and all three axes gives
\[
  \textup{matched-cell incidences}
  \ \ge\ \frac65|B_R|-6(2R+1).
\]
A matching supplies at most $|B_R|$ incidences.  At $R=20$, however,
\[
  \frac65(1261)-246=1267.2>1261,
\]
which is a contradiction.
\end{proof}

\begin{figure}[tbp]
  \centering
  \begin{tikzpicture}[scale=.55]
    \HexPatch{0/0,1/0,2/0,3/0,4/0,5/0,6/0,7/0,8/0,9/0,10/0,11/0}
    \HexMark{5}{0}{dot}{}
    \HexMark{7}{0}{dot}{}
    \draw[distance annotation]
      ({sqrt(3)*5},1.35) -- node[midway,fill=white,inner sep=1.5pt,
        font=\scriptsize] {$\delta=2$} ({sqrt(3)*7},1.35);
    \foreach \s/\y in {2/-1.55,3/-1.95,4/-2.35,5/-2.75}{
      \pgfmathsetmacro{\e}{\s+5}
      \draw[line width=.55pt]
        ({sqrt(3)*\s},\y+.18) -- ({sqrt(3)*\s},\y)
        -- ({sqrt(3)*\e},\y) -- ({sqrt(3)*\e},\y+.18);
    }
  \end{tikzpicture}
  \caption{A pair at axis-distance $\delta=2$ is contained in exactly
  $6-\delta=4$ length-six windows.  The four brackets show their possible
  starts.}
  \label{fig:pairing-density}
\end{figure}

The same density calculation applies to other translation-invariant line
systems with negligible ball boundary.

\begin{corollary}[Translation-invariant line grids; source T8.1]
\label{cor:T8.1}
Consider $k$-in-a-row on a point lattice with $g$ pairwise nonparallel,
translation-invariant line directions.  Suppose every lattice line has a
length-$k$ window at every offset and the lattice balls have
boundary-to-volume ratio tending to zero, as in the F{\o}lner property of
$\mathbb Z^2$.  A pairing requires per-line matched density $2/(k-1)$ and
hence
\[
  g\frac{2}{k-1}\le1
  \qquad\Longrightarrow\qquad k\ge2g+1.
\]
For the square grid, $g=4$ gives $k\ge9$, consistent with the classical
existence of pairing draws for $9$-in-a-row.  For the hex grid, $g=3$ gives
$k\ge7$, so $6$-in-a-row misses the bound by one.  Existence at the threshold
requires a separate construction and is not claimed.  The conclusion applies
to lattices satisfying the stated hypotheses, not to arbitrary grids.
\end{corollary}

The global-pairing approach is therefore closed.  Partial pairings on
constrained regions are not ruled out, but they cannot be position-independent.
